% AIDC 2026 body — a DISTINCT short position/WIP paper for the AIDC workshop
% (Agentic AI in Offensive and Defensive Cyber Operations). This is intentionally
% NOT the SSR SoK: the thesis here is sorting agent-authorization threats by whether
% a defender can test them adversarially before deployment, plus a paired generator
% and analyzer for the testable subset. Class-agnostic; the wrapper sets
% class/title/authors and the bibliography.
\ifdefined\ifsubmission\else\expandafter\newif\csname ifsubmission\endcsname\submissiontrue\fi
\providecommand{\gate}[1]{\ifsubmission #1\else \textbf{[PENDING: #1]}\fi}
\providecommand{\draftonly}[1]{\ifsubmission\else \textbf{[NOTE: #1]}\fi}
\providecommand{\yes}{\ensuremath{\surd}}
\providecommand{\no}{\ensuremath{\times}}
\emergencystretch=3em

\begin{abstract}
An autonomous agent acts on a bearer token. If an attacker can steer that agent
through untrusted input, such as a poisoned document or a malicious tool
description, the attacker gains whatever the token permits. Authorization is
therefore part of an agent's attack surface, and it has one useful property that the
model's behavior does not: a defender can test it adversarially before deployment.
This paper takes that offensive view and uses it to sort the defensive problem. We
place fifteen documented agent-token threats into two groups by a single question:
can a defender reproduce the attack offline, from one token or server manifest? Six
can be reproduced this way. The other nine cannot, because they depend on a live
runtime, such as a consent flow, a human approver, or a workload-identity layer that
most deployments do not yet run. For the six, we release an open-source generator
that mints adversarial agent tokens and an analyzer that rejects them, with every
finding cited to the OAuth RFCs, the MCP authorization specification, or the OWASP
Agentic Top~10. Both run in CI or next to the agent, so a team can red-team its own
authorization before shipping. The nine that cannot be tested offline set out a
runtime-defense agenda. This is work in progress: the instrument is implemented, and
a prevalence measurement across public deployments is ongoing and is not reported
here.
\end{abstract}

\section{Introduction}
A tool-using agent presents a bearer token on every call it makes for a user, either
an OAuth access token or an MCP authorization token. The agent is autonomous, and its
inputs are not. It reads web pages, summarizes documents, and follows tool
descriptions, any of which an attacker may control. Prompt injection does not need a
memory-safety bug to cause damage. It needs a token that grants more than the current
action requires.

The resulting failures are ordinary OAuth failures, and the agent's autonomy makes
them easier to reach~\cite{rfc9700,mcpauth}. A token that is not bound to the server
it is sent to can be replayed against a different server, or passed through it. An
on-behalf-of token with no verifiable delegation chain leaves the question ``who
authorized this?'' unanswerable. A wildcard scope such as \texttt{admin:*} gives an
injected agent far more access than any one task needs. A token with no expiry stays
useful to whoever obtains it.

Much of this is testable before an agent ever meets an adversary. A defender can mint
the same malicious tokens an attacker would forge, then confirm that its own verifier
rejects them. That is an inexpensive red-team of the authorization layer, and it can
run on every build. It also has a clear limit. Some agent-authorization attacks
cannot be decided from a token at all and appear only against a live deployment.
Separating the two cases is the work of this paper.

\noindent\textbf{Contributions.} First, we sort the agent-authorization threat
surface by offensive testability: what a defender can reproduce before deployment,
against what only a runtime defense can stop. Second, we release an open-source pair
for the testable subset, a generator of adversarial agent tokens and a stateless
analyzer that rejects them with cited findings. Third, we argue that the threats a
defender cannot test offline are a runtime-defense agenda, not a gap in tooling. This
is a position and work-in-progress paper. It reports a method and an instrument, not
a measurement.

\section{The agent-authorization attack surface}\label{sec:surface}
Table~\ref{tab:tax} lists fifteen threats in six categories. Each is tied to its
normative source and to the invariant an attacker breaks. The final column records
the distinction this paper is built on: whether a defender can reproduce the attack
from a single token or server manifest, which we call pre-deployment testable, or
whether the attack only appears against a live, stateful deployment.

\begin{table}[t]
\centering\footnotesize
\caption{Agent-authorization threats by category. ``Pre-dep.'': a defender can mint
the malicious input and check for rejection offline, from one token or manifest.
Otherwise the attack needs a live runtime (a consent flow, a human approver, or a
workload-identity substrate) and belongs to runtime defense.}
\label{tab:tax}
\begin{tabular}{@{}llp{3.7cm}c@{}}
\toprule
ID & Sev. & Attack / broken invariant & Pre-dep.? \\
\midrule
\multicolumn{4}{@{}l}{\emph{Audience binding} (RFC 8707/9728~\cite{rfc8707,rfc9728})}\\
MCP-01 & crit & token passthrough / wrong-audience accepted & \yes \\
MCP-02 & high & missing resource indicator / unbound audience & \yes \\
\multicolumn{4}{@{}l}{\emph{Delegation} (RFC 8693~\cite{rfc8693})}\\
DEL-01 & high & missing or forged \texttt{act} claim & \yes \\
DEL-02 & high & delegation-chain / transitive-trust abuse & \yes \\
DEL-03 & high & \texttt{may\_act} not enforced on exchange & \no \\
DEL-04 & med  & impersonation where delegation is required & \no \\
\multicolumn{4}{@{}l}{\emph{Scope}}\\
SCOPE-01 & high & over-scoped agent credential & \yes \\
SCOPE-02 & high & non-expiring / long-lived token & \yes \\
SCOPE-03 & high & agent reuses the human's root credential & \no \\
\multicolumn{4}{@{}l}{\emph{Agent identity} (WIMSE/SPIFFE~\cite{wimse,spiffe})}\\
AID-01 & high & session used as authentication & \no \\
AID-02 & med  & no distinct workload/agent identity & \no \\
\multicolumn{4}{@{}l}{\emph{Confused deputy} (CWE-441~\cite{cwe441,hardy1988confused})}\\
CD-01 & crit & confused deputy via static client ID & \no \\
CD-02 & high & dynamic client registration abuse (RFC 7591~\cite{rfc7591}) & \no \\
\multicolumn{4}{@{}l}{\emph{Agency} (OWASP~\cite{owaspllm,owaspagentic})}\\
AGT-01 & crit & prompt-injection-driven privilege escalation & \no \\
AGT-02 & high & tool poisoning / rug-pull & \no \\
\bottomrule
\end{tabular}
\end{table}

Six of the fifteen are pre-deployment testable: MCP-01, MCP-02, DEL-01, DEL-02,
SCOPE-01, and SCOPE-02. Each is a property of a token's own claims, read against a
stated policy. The audience it names, the \texttt{act} chain it carries, the scopes
it requests, and the lifetime it declares are all present in the token, so an attacker
forges them offline and a defender reproduces them the same way. The remaining nine
are not testable this way. A confused deputy (CD-01, CD-02) is a property of a consent
flow and of how a proxy matches \texttt{redirect\_uri}~\cite{mcpauth,rfc7591}.
Prompt-injection escalation and tool poisoning (AGT-01, AGT-02) are properties of a
running agent and the tools it has been given~\cite{owaspllm,owaspagentic}. Identity
confusion (AID-01, AID-02) assumes a workload-identity layer~\cite{wimse,spiffe} that
is usually not deployed, so there is nothing to check a token against.

\section{Offense: generating adversarial agent tokens}\label{sec:offense}
For each testable class we generate the token an attacker would present. Generation is
direct because the invariant is explicit. For audience binding, we emit a token whose
resource indicator names a different server than the one it is sent to, or names none
(MCP-01, MCP-02). For delegation, we emit an on-behalf-of token with a missing,
malformed, or over-deep \texttt{act} chain (DEL-01, DEL-02). For scope, we emit a
credential that carries a wildcard or an otherwise over-broad scope beyond a declared
allowlist (SCOPE-01). For expiry, we emit a token with no \texttt{exp}, or a lifetime
past a declared maximum (SCOPE-02).

These are the inputs a real attacker crafts once it can influence what an agent
requests, or can replay a token the agent was issued. Generating them on purpose is
the offensive half of a pre-deployment red-team: the defender produces the
adversary's tokens and asks whether its own verifier lets any through. A token that is
accepted is a confirmed weakness that an attacker could exploit, with no further
analysis needed to interpret it. Figure~\ref{fig:loop} shows one such token and the
findings it draws.

\begin{figure}[t]
\footnotesize
\begin{verbatim}
adversarial agent token (claims)
  iss   = https://issuer.example
  sub   = agent-7
  scope = "admin:* files:read"
  aud   = (absent)   exp = (absent)

analyzer output (each finding cited)
  MCP-02   high  no audience; token is
           not bound to a resource  [RFC 8707]
  DEL-01   high  act claim missing;
           delegation unverifiable  [RFC 8693]
  SCOPE-01 high  scope "admin:*" is
           over-broad          [OWASP Agentic]
  SCOPE-02 high  no exp; a non-expiring
           credential               [RFC 8725]
\end{verbatim}
\caption{One adversarial agent token and the analyzer's findings. The generator emits
the token; the analyzer rejects it and cites a source for each finding. Run as a CI
gate, any token the verifier accepts fails the build.}
\label{fig:loop}
\end{figure}

\section{Defense: a deployable analyzer}\label{sec:defense}
The defensive half is a stateless analyzer. It decodes a presented agent token and
reports findings against a policy, each citing a normative source: audience binding
(MCP-01, MCP-02; RFC 8707/9728 and the MCP authorization specification), on-behalf-of
delegation (DEL-01, DEL-02; RFC 8693), scope minimization (SCOPE-01; OWASP Agentic),
and expiry (SCOPE-02; RFC 8725). It runs on one token, makes no network calls, and so
fits either a CI check against fixtures or a client-side check next to the agent.
Nothing leaves the operator.

Paired with the generator of Section~\ref{sec:offense}, the analyzer closes a loop a
team can run on every build. Mint the adversarial tokens, confirm the verifier rejects
each one, and fail the build if any is accepted. We release both components as open
source under a permissive license. The repository link is withheld here for anonymous
review.

\section{Where defense has to live}\label{sec:line}
The split of six against nine is a claim about where a defense has to run, not only
about what a tool happens to cover. Classic OAuth machinery, including resource
indicators, token exchange, and protected-resource metadata, does model several of
the nine runtime threats~\cite{rfc8707,rfc8693,rfc9728}. Modeling is not verifying,
and common MCP deployments rarely enforce these mechanisms, so the threats fall to the
deployer, which is where they are usually missed.

The nine resolve into a concrete defense agenda. CD-01 and CD-02 need per-client
consent and exact \texttt{redirect\_uri} matching at the proxy. DEL-03, DEL-04, and
SCOPE-03 need the server to enforce \texttt{may\_act} on token exchange and to keep
the agent's credential distinct from the human's. AID-01 and AID-02 need a deployed
workload-identity substrate~\cite{wimse,spiffe}. AGT-01 and AGT-02 need per-action
authorization, or a human in the loop, so that an injection cannot escalate on its
own~\cite{owaspagentic}. Every one of these is a runtime property, and none can be
settled by reading a token. Knowing which side of the line an attack sits on tells a
team what to automate in CI today and what to treat as an open runtime problem.

\section{Status and limitations}\label{sec:status}
This is work in progress. The generator and the analyzer are implemented for the six
pre-deployment-testable classes and released, and the loop between them runs in CI and
client-side. We do not report a prevalence measurement. A study of how often these
weaknesses occur across public MCP servers and reference agent-token deployments is
ongoing, with a hand-validated sample, and we will report it, including its skew
toward early-adopter deployments, rather than imply it here. Two limitations bound the
method. The MCP specification is still changing, so the taxonomy is a dated snapshot
that will need revision. The testable line depends on what a token or manifest
exposes: a deployment that hides relevant state reduces what any offline red-team can
conclude.

\section{Related work}\label{sec:related}
Our building blocks are the OAuth~2.0 framework~\cite{rfc6749} and its security best
practice~\cite{rfc9700}, token exchange~\cite{rfc8693}, resource
indicators~\cite{rfc8707}, protected-resource metadata~\cite{rfc9728}, dynamic client
registration~\cite{rfc7591}, JWT best practices~\cite{rfc8725}, and the MCP
authorization specification~\cite{mcpauth}. The confused-deputy problem is
longstanding~\cite{hardy1988confused,cwe441}. Practitioner taxonomies catalog
agentic risk~\cite{owaspllm,owaspagentic}, and recent drafts address on-behalf-of
agent authentication~\cite{oboagents} and workload identity~\cite{wimse,spiffe}.
Offensive work on agents has so far concentrated on prompt injection and model
manipulation. We instead treat the authorization layer as the thing to attack and
defend, and we contribute both the sort by pre-deployment testability and a paired
instrument for the testable subset.

\section*{Ethics considerations}
The work uses only public specifications, advisories, and CVEs. The generator produces
adversarial tokens for a defender to fire at its own verifier. It is deny-by-default
and documented for staging rather than production use. We probe no third-party live
agents or MCP servers and use no real user tokens. The analyzer runs on tokens the
operator already holds and sends nothing anywhere. Any future measurement will
characterize population prevalence rather than single deployments, and will follow
coordinated disclosure for any identifiable high-severity finding. No personal data or
human subjects are involved.

\section*{Open science}
We release the generator, the analyzer, and the machine-readable threat taxonomy, with
each threat recorded alongside its invariant and normative sources, under an open
license. The offense-and-defense loop is therefore reproducible and the taxonomy is
open to inspection. \gate{A DOI-pinned artifact will accompany the camera-ready
version.}

\section{Conclusion}
Agent authorization is both reachable by attackers and, in large part, testable by
defenders. Six of fifteen documented agent-token threats can be exercised
adversarially before deployment. We released a generator and an analyzer that let a
team run that red-team in CI, and we argued that the nine threats which cannot be
tested offline are the runtime-defense agenda for agentic systems.
